Las siguientes líneas de trabajo futuro y de operación se formulan precisando, en cada caso, qué debe hacerse, mediante qué enfoque y con qué recursos.

\begin{enumerate}
	\item Modelos basados en árboles. Queda pendiente implementar y evaluar los árboles de decisión (contemplados inicialmente y no desarrollados) junto con sus extensiones de conjunto, como los bosques aleatorios y los árboles potenciados por gradiente. Para que la comparación con la red neuronal sea directa, deben entrenarse sobre las mismas dos formas y la misma partición temporal (cohortes 2010–2016 para entrenamiento, 2017 para validación y 2018 para prueba), evaluarse con las mismas métricas (MAE, RMSE y R²) y ajustar sus hiperparámetros mediante búsqueda sobre el conjunto de validación. Esto puede realizarse con las bibliotecas Scikit-learn, para los árboles y bosques, y XGBoost o LightGBM, para el potenciado por gradiente, reutilizando el conjunto emparejado ya construido; al no requerir estandarización previa, su incorporación al flujo existente es directa.
	\item Variables del proceso universitario. Dado que la capacidad explicativa está acotada por la naturaleza del problema, el siguiente paso es incorporar variables propias de la trayectoria en la educación superior como el promedio acumulado, el número de asignaturas reprobadas o la permanencia, que podrían explicar una porción mayor de la varianza atribuible al valor agregado. Se propone integrar estas variables al conjunto de datos emparejado mediante el identificador de cada estudiante e incorporarlas como una tercera configuración de entrada para evaluar si aportan mejoras en la capacidad explicativa de los modelos frente a las configuraciones actuales. Para hacerlo, sería necesario contar con acceso a los sistemas de información académica de la Universidad, especialmente a los relacionados con registro y control académico, además de definir acuerdos para el uso de la información y aplicar mecanismos de anonimización que garanticen la protección de los datos personales.
	\item Cálculo ponderado del puntaje global de Saber Pro. También se recomienda evaluar una versión del puntaje global de Saber Pro en la que la competencia de Inglés tenga una ponderación menor, siguiendo un criterio similar al utilizado por el ICFES en la prueba Saber 11. Esta modificación puede implementarse dentro del procedimiento de cálculo del puntaje global que ya está incorporado en el sistema. Posteriormente, los resultados podrían compararse con los obtenidos mediante el promedio simple para analizar el impacto que esta decisión tiene sobre la estimación del valor agregado institucional. Una ventaja de esta propuesta es que puede desarrollarse utilizando la información disponible actualmente, sin requerir nuevas fuentes de datos.
	\item Modelado diferenciado por programa y cohorte. Conviene avanzar de un comportamiento institucional agregado hacia un análisis por programa académico y por cohorte, que permita identificar fortalezas y debilidades específicas. El enfoque consiste en estratificar el entrenamiento y la evaluación por subgrupo, o incorporar el programa como variable, reportando el valor agregado por programa siempre que el volumen de cada subgrupo lo permita. Para ello se dispone del dato de programa presente en los registros y del propio tablero, cuyo filtro actual por institución puede extenderse a un filtro por programa.
	\item Efecto de la pandemia. Es necesario estudiar de manera dedicada el marcado deterioro del valor agregado observado a partir de 2019. Se sugiere establecer líneas base ajustadas a ese periodo o introducir un indicador de periodo pandémico, y aplicar técnicas de detección de cambio estructural para separar el efecto coyuntural del institucional. Este análisis puede realizarse sobre los datos ya disponibles entre 2016 y 2023, sin recolección adicional.
	\item Automatización de la carga de nuevos datos. El sistema opera actualmente solo con los datos existentes: aunque están desarrollados los procesos que depuran la información, ejecutan la extracción, transformación y carga hacia PostgreSQL y entrenan las redes, no se dispone de una opción sencilla, de cara al usuario, para incorporar nuevos conjuntos de datos. Además, las redes neuronales se encuentran actualmente configuradas de acuerdo con las cohortes de entrenamiento y prueba utilizadas en este estudio, y gran parte del flujo de procesamiento depende de parámetros definidos manualmente para los años y conjuntos de datos disponibles. Por esta razón, una línea de mejora futura consiste en facilitar la incorporación de nueva información, de manera que el sistema pueda adaptarse a futuras cohortes sin requerir modificaciones directas en el código. Para lograrlo, sería conveniente automatizar elementos que actualmente permanecen definidos de forma fija, como los años analizados, la definición de los conjuntos de entrenamiento, validación y prueba, y las rutas de acceso a las tablas de datos. La idea es que el propio sistema pueda identificar esta información a partir de los datos disponibles y ajustar su configuración de manera automática. Adicionalmente, se propone incorporar una interfaz que permita cargar nuevos archivos de forma sencilla y segura. Esta funcionalidad debería verificar la estructura y el formato de los datos, aplicar los procesos de limpieza y normalización ya establecidos y ejecutar de manera controlada tanto la integración de la información como el reentrenamiento de los modelos cuando sea necesario. Este proceso puede aprovechar y reutilizar los procesos de depuración, ETL y entrenamiento ya existentes, encapsulados tras la interfaz del propio tablero o un módulo de administración, y apoyándose en la infraestructura desplegada (PostgreSQL y el servidor de la aplicación) para orquestar el proceso.
	\item Operación y mantenimiento del sistema. Para garantizar la continuidad y vigencia del sistema, se recomienda reentrenarlo a medida que nuevas cohortes completen ambas pruebas, conservar como artefacto canónico la predicción promediada sobre múltiples semillas y asegurar que tanto el cálculo del global como el tablero consuman ese artefacto del ensemble. Esto se sostiene sobre la infraestructura ya desplegada —PostgreSQL, Gunicorn, Nginx y systemd— y sobre el manual técnico de despliegue, mantenido conforme a las normas ISO aplicables.
\end{enumerate}